% NOTES ONLY, to be written. Trimmed 2026-11-07 from thesis/methodology/03_prediction_model.tex
% (full draft in git, commit e4eb857).
% Design: root CLAUDE.md "Decided". Settings come from CGPS_data/progression_who_pooled.py, which is
% not in CGPS_data/thesis_final/ and still runs on the Medis trees with a progression outcome; the
% text must follow the clean-cohort version once it exists there.
%
% The section should contain:
% - One segment-level complementary log-log GLMM over the nine segments, random intercept per
%   participant (sec:tech-segment-glmm). One row per segment, outcome incident coronary plaque in
%   that segment (sec:method-population).
% - Segment risks pooled into one person risk of plaque in any of the nine segments
%   (sec:tech-aggregation); the person risk is what is evaluated.
% - Fit: maximum likelihood, Gauss-Hermite quadrature over the random intercept. Open: node count
%   (30 in the script, checked by a refit with doubled nodes).
% - Open (added 2026-11-13): random-effects structure. The person intercept makes every pair of
%   segments equally correlated (pLAD-mLAD as pLAD-dRCA); segments on one artery likely share more
%   (inflow, plaque across segment boundaries). Effect: LR test rejects too often (already in
%   sec:tech-glmm-likelihood), and the person risk is overstated if plaques cluster within an artery.
%   Options: (a) add a person x artery intercept (LAD, LCx, RCA), nested in person; about three
%   segments per artery and few events, so its variance may not be identifiable; (b) check residual
%   correlation for same-artery vs different-artery pairs after the fit; (c) if (a) fails, keep the
%   person intercept and state it as a limitation in the Discussion. Intervals from person-level
%   resampling (CV, bootstrap) hold whatever the dependence; only the LR test leans on the intercept.
%   Not a case for GEE: an exchangeable working correlation makes the same assumption, and GEE gives
%   no joint distribution for the person risk. Theory: sec:tech-nested.
% - Base model: segment position (pLAD reference), log segment length, dominance, log follow-up
%   time, SCORE2 risk factors. Open: follow-up as term or offset; recalibrated SCORE2
%   (sec:tech-recalibration) vs refitted SCORE2 factors, compared in this pooled model.
% - Extended model: base + geometry block (four features, sec:method-selection: curvature, taper,
%   bifurcation angle, daughter asymmetry ln(D_LCx/D_LAD); each its own linear term; Finet's ratio
%   dropped 2026-11-13, see 05c header),
%   standardised within segment type on the training participants. Open: which segment rows carry the
%   left main angle and asymmetry (script: pLAD and pLCx rows only).
% - Secondary (decided 2026-11-11): angle x asymmetry, one prespecified term, estimated and reported,
%   outside the LR test of the block (a product needs its main effects; support is one simulation
%   study on synthetic trees, garcha2025sensitivity). sec:method-ratio points here.
% - One prespecified LR test of the geometry block (eq:tech-lr), both models on all participants.
%   Delta C, calibration slope and intercept, Brier score and log loss estimated with intervals,
%   never tested (sec:tech-added-value).
% - Where: read descriptively from the same model (feature hazard ratios, segment intercepts,
%   within-participant C, secondary curvature/taper x artery test). No separate where model.
% - Subsection Cross-Validation (sec:method-cv): folds by participant, all nine segments in one
%   fold, stratified on developing plaque. Same folds for base and extended model; standardisation,
%   coefficients and random-intercept variance estimated on training participants of each fold.
%   Report mean over repeats with 2.5th/97.5th percentiles and the paired difference. Open: folds and
%   repeats (10 x 50 in the script); percentiles show split variation, not sampling uncertainty,
%   so state whether difference intervals come from a bootstrap of the whole procedure.
\section{Risk Model and Evaluation}
\label{sec:method-models}

\todonote{To be written: the pooled cloglog GLMM, its base and extended model, the LR test of the geometry block and the estimated performance measures. See the notes in the source.}

\subsection{Cross-Validation}
\label{sec:method-cv}

\todonote{To be written: person-level folds, stratified on the outcome, and how the measures are summarised.}
